\title{Trabalho Prático: Implementação e Análise - TCP vs. TLS}
\author{Prof. Gabriel Rodrigues Caldas de Aquino}
\date{Compilado em: \\ \today}

\begin{document}


\maketitle

\section{Objetivo}
Desenvolver um programa cliente em Python capaz de enviar um arquivo de texto para um servidor.
O envio deve ser realizado de duas formas: utilizando conexão normal (sem criptografia) e utilizando TLS para transmissão segura.


\textbf{Entrega: 05/12/2025}

\section{Requisitos do Trabalho}

\subsection{Implementação do Programa}
\begin{itemize}
    \item Criar um \textbf{cliente Python} que envie arquivos de texto para um servidor.
    \item Deve-se Implementar \textbf{duas formas de transmissão}:
          \begin{enumerate}
              \item \textbf{Sem TLS:} envio em texto plano.
              \item \textbf{Com TLS:} envio utilizando criptografia TLS.
          \end{enumerate}
\end{itemize}

\subsection*{2. Relatório}
O relatório deve conter:
\begin{itemize}
    \item Explicação do funcionamento do TLS, incluindo handshake e criptografia de dados.
    \item Comparação entre a implementação sem TLS e com TLS.
    \item Captura de pacotes de rede mostrando a diferença entre envio cifrado e envio em texto plano (pode ser feito com Wireshark ou ferramenta equivalente).
    \item Análise do \textbf{overhead} introduzido pelo uso de TLS (ex.: tamanho do pacote, tempo de envio).
    \item Discussão sobre o overhead.

\end{itemize}

\section*{Entrega}
\begin{itemize}
    \item Código-fonte do cliente e servidor (Python) nas implementações com e sem TLS.
    \item Relatório em PDF ou documento de texto estruturado com todas as análises solicitadas.
\end{itemize}

\section*{Avaliação}
\begin{itemize}
    \item Funcionalidade do cliente e servidor em ambos os modos.
    \item Clareza e detalhamento do relatório, incluindo evidências das capturas de pacotes.
    \item Correção e profundidade na análise do overhead e do funcionamento do TLS.
\end{itemize}

\end{document}